% BHU PYQ -- Professional Exam Template
\documentclass[11pt,a4paper]{article}

\usepackage[a4paper,top=2.5cm,bottom=2.5cm,left=2.8cm,right=2.8cm,headheight=14pt]{geometry}
\usepackage{amsmath,amssymb}
\usepackage[T1]{fontenc}
\usepackage[utf8]{inputenc}
\usepackage{lmodern}
\usepackage{microtype}
\usepackage{fancyhdr}
\usepackage{enumitem}
\usepackage{hyperref}
\usepackage{lastpage}
\usepackage{parskip}

\hypersetup{colorlinks=true,urlcolor=black,linkcolor=black,
  pdftitle={B.Sc. (Hons.) Zoology Sem-V ZOB-503 -- BHU PYQ}}

\pagestyle{fancy}
\fancyhf{}
\renewcommand{\headrulewidth}{0.4pt}
\renewcommand{\footrulewidth}{0.4pt}
\fancyhead[L]{\small\textit{Banaras Hindu University}}
\fancyhead[R]{\small\textit{Zoology Paper: ZOB-503}}
\fancyfoot[C]{\small Page \thepage\ of \pageref{LastPage}}

\newlist{parts}{enumerate}{3}
\setlist[parts,1]{label=(\alph*),leftmargin=*,itemsep=4pt,topsep=3pt}
\setlist[parts,2]{label=(\roman*),leftmargin=*,itemsep=2pt,topsep=2pt}
\setlist[parts,3]{label=(\arabic*),leftmargin=*,itemsep=2pt,topsep=2pt}
\newcommand{\pts}[1]{\hfill\textit{[#1]}}

\begin{document}

\begin{center}
  {\large\bfseries BANARAS HINDU UNIVERSITY}\\[2pt]
  {\normalsize B.Sc. (Hons.) Semester V Examination, 2016--17}\\[6pt]
  \rule{0.8\linewidth}{0.5pt}\\[6pt]
  {\large\bfseries Zoology}\\[4pt]
  {\large\bfseries Paper No. ZOB-503 : Biochemistry and Molecular Biology}\\[4pt]
  {\normalsize Time: 3 hours \hfill Full Marks: 70}
\end{center}

\rule{\linewidth}{0.4pt}

\smallskip
\noindent\textit{Note: Write your Roll No. at the top immediately on the receipt of this question paper. The figures in the right-hand margin indicate marks. Answer six questions including Question 1 which is compulsory. Coloured illustrations permitted.}

\bigskip
\noindent\textbf{Question 1.} Answer the following parts:
\begin{parts}
  \item State whether the following statements are True or False with brief justification:\pts{$8 \times 2 = 16$}
  \begin{parts}
    \item B-DNA and A-DNA are covalently closed circular DNA.
    \item CPSF inhibits capping of RNA.
    \item SeqA binds to RNA to initiate transcription.
    \item $(\text{pK}_1 + \text{pK}_2)/2$ is a common formulae to calculate pI of all types of $\alpha$-amino acids.
    \item A disulphide bond in a protein is generated between a cys and a met residue.
    \item Carnitine shuttle aids in oxidative phosphorylation.
    \item A sigmoidal kinetics is shown by a multimeric enzyme.
    \item Cerebrosides are derivatives of sphingolipids.
  \end{parts}

  \item Match each item from Column--X with one from Column--Y:\pts{$4 \times \frac{1}{2} = 2$}
  \begin{center}
  \begin{tabular}{llp{1.5cm}ll}
    \hline
    \multicolumn{2}{c}{\textbf{Column--X}} & & \multicolumn{2}{c}{\textbf{Column--Y}} \\
    \hline
    (i) & Palindrome & & (1) & Tryptophan \\
    (ii) & UAA & & (2) & Restriction endonucleases \\
    (iii) & UGA & & (3) & Termination \\
    (iv) & GUG & & (4) & Initiation of translation \\
    & & & (5) & RNA polymerase-I \\
    \hline
  \end{tabular}
  \end{center}
  \medskip

  \item Define the following:\pts{$2 \times 1 = 2$}
  \begin{parts}
    \item Codon
    \item Protein domain.
  \end{parts}
\end{parts}

\pagebreak

\medskip\noindent\textbf{Question 2.} Answer the following parts:
\begin{parts}
  \item Represent steps of one cycle of Edman's reaction with a polypeptide chain.\pts{5}
  \item Using Lineweaver-Burk equation, justify all the kinetic parameters shown in the LB plot.\pts{5}
\end{parts}

\medskip\noindent\textbf{Question 3.} Answer the following parts:
\begin{parts}
  \item Compare stabilization of $\alpha$-helix and $\beta$-pleated sheet structures of polypeptides.\pts{6}
  \item Write pay-off reactions of the glycolytic pathway.\pts{4}
\end{parts}

\medskip\noindent\textbf{Question 4.} Answer the following parts:
\begin{parts}
  \item Explain chemiosmotic concept of ATP synthesis.\pts{5}
  \item Write de-carboxylation reactions of the TCA cycle.\pts{5}
\end{parts}

\medskip\noindent\textbf{Question 5.} With the help of well-labeled diagrams, explain the following:\pts{$2 \times 5 = 10$}
\begin{parts}
  \item Formation of pre-initiation complex at promoter for RNA Pol II
  \item Mechanism of termination of translation in prokaryotes.
\end{parts}

\medskip\noindent\textbf{Question 6.} Describe compositions and functions of the following:\pts{$2 \times 5 = 10$}
\begin{parts}
  \item Telomerase
  \item Nucleosome.
\end{parts}

\medskip\noindent\textbf{Question 7.} Differentiate between the following:\pts{$2 \times 5 = 10$}
\begin{parts}
  \item Cloning and Expression vectors
  \item Primosome and Replisome.
\end{parts}

\medskip\noindent\textbf{Question 8.} Write short notes on the following:\pts{$2 \times 5 = 10$}
\begin{parts}
  \item Phospholipids
  \item Ribosome.
\end{parts}

\vfill
\rule{\linewidth}{0.4pt}
\begin{center}\small
\href{https://bhu-exam.vercel.app/}{Click Here}\\[4pt]\href{https://me-aryan.vercel.app/}{Click Here}\\[4pt]\href{https://quashan-me.vercel.app/}{Click Here}
\end{center}

\end{document}
